\section{SSI 的定位与策略}

\subsection{Straight-Shot Super Intelligence？}

SSI 最初的计划是"直奔超级智能"（straight-shot superintelligence），跳过市场竞争。Ilya 列出了正反两方面的考量：

\textbf{支持直奔}：
\begin{itemize}
    \item 不受日常市场竞争的困扰，可以专注研究
    \item 避免被迫做出困难的 trade-off
\end{itemize}

\textbf{反对直奔}：
\begin{itemize}
    \item 让世界看到强大 AI 是有价值的——"communicate the AI, not the idea"
    \item 如果时间线很长，纯研究模式不可持续
    \item 最好的 AI 被使用并在使用中改进（类比飞机安全的提升来自实际飞行）
\end{itemize}

\subsection{Ideas 的稀缺}

\begin{importantbox}{If ideas are so cheap, how come no one's having any ideas?}
Scaling 时代的遗产之一是：公司多于 idea。所有人在做同样的事情（RL on pre-trained models），真正的差异化 idea 稀缺。SSI 声称拥有不同的技术路径（围绕泛化和理解），如果这些 idea 被证明正确，将具有重要价值。
\end{importantbox}

\subsection{共同创始人离开}

Ilya 简要回应了联合创始人离开加入 Meta 的事件：SSI 当时正在以 320 亿美元估值融资，Meta 提出收购邀约。Ilya 拒绝了，但前联合创始人选择了加入 Meta（也是唯一从 SSI 加入 Meta 的人）。

\subsection{本章小结}

SSI 定位为"research 时代的公司"，拥有围绕泛化的差异化技术路径。原先的"直奔超级智能"策略可能因时间线和部署价值的考量而调整。


